% latexmk -pdf -pvc blatt13.tex
\documentclass{hott-übung}

\begin{document}

\setcounter{blattnummer}{13}
\setcounter{aufgabennummer}{0}

\blatt

\aufgabe{}
Berechne die Kohomologie \(H^1(S^1)\).

\aufgabe{}
Sei \(A\) ein Typ und \(B : A \to \mU\) ein abhängiger Typ über \(A\).
Zeige:
\begin{enumerate}
  \item \(\|\sum_{x:A}B(x)\|_n \simeq \|\sum_{x:A}\|B(x)\|_n\|_n\). {\tiny Vergleiche mit Blatt 9 A3.}
  \item Ist \(B(x)\) für jedes \(x : A\) ein \(n\)-zusammenhängender Typ, so ist \(\sum_{x : A} B(x)\) genau dann \(n\)-zusammenhängend, wenn \(A\) es ist.
\end{enumerate}
Seien nun Abbildungen \(f : X \to Y\), \(g : Y \to Z\) gegeben.
Folgere:
\begin{enumerate}[resume]
  \item Ist \(f\) \(n\)-zusammenhängend, so ist \(g\circ f\) genau dann \(n\)-zusammenhängend, wenn \(g\) es ist.
\end{enumerate}

\aufgabe{}
Sei \(f:A\to B\) \(n\)-zusammenhängend.
Zeige, dass \(\|f\|_n:\|A\|_n\to \|B\|_n\) eine Äquivalenz ist.

\aufgabe{}
Sei \(G\) eine abelsche Gruppe und \(K(G,2)\) ein punktierter \(1\)-zusammenhängender \(2\)-Typ, sodass \(\Omega^2 K(G,2) \simeq G\) ist.
Zeige, dass \(K(G,3) \colonequiv \|\Sigma K(G,2)\|_3\) ein punktierter \(2\)-zusammenhängender \(3\)-Typ ist, welcher \(\Omega^3 K(G,3) \simeq G\) erfüllt.

\end{document}
